% DriveMotorModule
\section{【応用】DriveMotorModule — PWM 2ピン制御 + 統合部品}

\begin{patternbox}{DriveMotorModule で学べるパターン}
\begin{itemize}
    \item Hブリッジによる方向制御（IN1/IN2 + PWM）
    \item 逆転保護（ブレーキ自動挿入）
    \item drive()メソッドによるSystemData非経由の直接制御
    \item 統合モジュールの「内部部品」としての使い方
\end{itemize}
\end{patternbox}

\subsection{2つの使い方}

\texttt{DriveMotorModule}は2つのモードで使える。

\begin{enumerate}
    \item \textbf{単体使用} — \texttt{outputModules[]}に登録し、\texttt{SystemData}経由で制御
    \item \textbf{統合モジュールの内部部品} — \texttt{ChassisModule}等が内部で\texttt{drive()}を呼ぶ
\end{enumerate}

\begin{lstlisting}[style=cppstyle, caption={DriveMotorModule.h — drive()メソッド}]
class DriveMotorModule : public IModule {
public:
    // SystemData経由の制御（単体使用時）
    void update(SystemData& data) override;

    // SystemData非経由の直接制御（統合モジュールから使用）
    void drive(float power);
    void brake();
};
\end{lstlisting}

\begin{pointbox}[update()とdrive()の使い分け]
\texttt{update()}はSystemDataの\texttt{DriveMotorData.power}を読んで\texttt{drive()}を呼ぶラッパー。
統合モジュールの内部部品として使う場合は\texttt{drive()}を直接呼び、
\texttt{update()}は空実装（呼ばれない）となる。
この設計により、同じモジュールコードを単体でも統合でも使い回せる。
\end{pointbox}

\subsection{逆転保護}

\begin{lstlisting}[style=cppstyle, caption={DriveMotorModule.cpp — 逆転保護}]
void DriveMotorModule::drive(float power) {
    Direction newDir = (power > 0) ? FORWARD
                     : (power < 0) ? REVERSE : STOP;

    // 正転<->逆転の直接切り替え時、1フレーム分のブレーキを挿入
    if ((_lastDir == FORWARD && newDir == REVERSE) ||
        (_lastDir == REVERSE && newDir == FORWARD)) {
        brake();
        _lastDir = STOP;
        return;  // 次のupdate()で実際の方向転換が行われる
    }
    // ... 実際のPWM制御 ...
}
\end{lstlisting}

\begin{pointbox}[1フレーム遅延によるモーター保護]
正転→逆転の直接切り替えはモーターやドライバICに大きな負荷がかかる。
1フレーム分（1ループ分）のブレーキを挿入することで、物理的に安全な方向転換を実現する。
3フェーズモデルの「毎ループ1回だけ呼ばれる」性質を利用した保護パターン。
\end{pointbox}
